\clearpage
\begin{summarybox}
    \textbf{\textcolor{CSummary}{一句话核心}}

    \textbf{外点引两条切线，连接两切点：将椭圆的 $x^2,y^2$ 分别换成 $x_0x,y_0y$，就得到切点弦所在直线。}
    \[
        \frac{x^2}{a^2}+\frac{y^2}{b^2}=1
        \quad\Longrightarrow\quad
        \boxed{\frac{x_0x}{a^2}+\frac{y_0y}{b^2}=1}.
    \]

    \medskip
    \textbf{\textcolor{CSummary}{先判定，再使用}}

    令 $S=\frac{x_0^2}{a^2}+\frac{y_0^2}{b^2}$：
    \begin{itemize}[leftmargin=*, itemsep=0.2em]
        \item $S>1$：$P$ 在外，两切点不同，公式表示切点弦。
        \item $S=1$：$P$ 在椭圆上，公式表示该点处的切线。
        \item $S<1$：$P$ 在内，不存在从 $P$ 引的两条实切线，不能称切点弦。
    \end{itemize}

    \medskip
    \textbf{\textcolor{CSummary}{三种证明，各有用途}}
    \begin{description}[style=nextline, leftmargin=2.8em, labelsep=0.5em, itemsep=0.25em]
        \item[\textbf{(1) 门槛最低：}] 椭圆缩放成圆 $\rightarrow$ 圆切线垂直半径 $\rightarrow$ 勾股定理产生乘积 $\rightarrow$ 换回椭圆。
        \item[\textbf{(2) 最快：}] 切点的切线过 $P$ $\rightarrow$ 两切点同满足乘积式 $\rightarrow$ 两点确定直线。
        \item[\textbf{(3) 常规通用：}] 设过 $P$ 与切点的直线 $\rightarrow$ 代入椭圆 $\rightarrow$ 相切时只有一个交点 $\rightarrow$ 一次项系数为零。
    \end{description}

    \medskip
    \textbf{\textcolor{CSummary}{高价值补充}}
    \begin{itemize}[leftmargin=*, itemsep=0.25em]
        \item 切点弦\textbf{不经过外点 $P$}：代入 $P$ 时左端为 $S>1$。
        \item 切点弦的中点为 $\boxed{M=(x_0/S,\ y_0/S)}$，所以 $O,M,P$ 共线。
        \item 已知切点弦 $Ax+By=C$（$C\ne0$），反求外点：
        $\boxed{x_0=a^2A/C,\ y_0=b^2B/C}$，并检查 $S>1$。
    \end{itemize}

    \textbf{考场检查：} 外点条件是否成立？写的是 $AB$ 还是切线 $PA$？分母是否仍是 $a^2,b^2$？
\end{summarybox}
